\documentclass[12pt,a4paper]{article}

\usepackage{fontawesome5}
\providecommand{\faPencilAlt}{\faIcon{pencil-alt}}

\usepackage{styles/couleurs}
\usepackage{styles/mise_en_page}
\usepackage{styles/boites}
\usepackage{styles/maths}
\usepackage{styles/figures}

\usepackage{enumitem}
\usepackage{array}
\usepackage{tabularx}
\usepackage{amsmath}
\usepackage{amssymb}
\usepackage{graphicx}
\usepackage{tikz}
\usepackage{xparse}

\usetikzlibrary{calc,arrows.meta}

\setlist[itemize]{leftmargin=1.4em,itemsep=0.15em,topsep=0.2em}
\setlist[enumerate]{leftmargin=*,itemsep=0.22em,topsep=0.25em}

\newcommand{\lignepointillee}{%
  \noindent\dotfill\par
  \vspace{0.45em}
}

\newcommand{\lignesreponse}[1]{%
  \begin{tcolorbox}[
    enhanced,
    colback=white,
    colframe=NeutreC!40,
    boxrule=0.5pt,
    arc=1.5mm,
    left=3mm,
    right=3mm,
    top=2mm,
    bottom=1mm,
    before skip=2pt,
    after skip=2pt
  ]
    \foreach \i in {1,...,#1}{\lignepointillee}
  \end{tcolorbox}
}

\newcommand{\imageactivite}[3]{%
  \IfFileExists{#3}{%
    \includegraphics[width=#1,height=#2,keepaspectratio]{#3}%
  }{%
    \fbox{%
      \parbox[c][#2][c]{#1}{%
        \centering\sffamily\small
        Image à intégrer\par
        \vspace{0.3em}
        \texttt{\detokenize{#3}}
      }%
    }%
  }%
}

\newcommand{\titreseq}[1]{%
\begin{center}
  {\sffamily\LARGE\bfseries\color{TitreC}
  Séquence 5 -- Statistiques et proportionnalité}\\[0.2em]
  {\sffamily\large\bfseries\color{SavoirC}#1}
\end{center}
}

\begin{document}

\titreseq{Progression professeur}

\remarque{Durée indicative : huit séances d'apprentissage de 50 à 55 minutes,
deux interrogations courtes et une évaluation finale.

La fiche d'exercices hors classe est distribuée en début de séquence.
Les exercices à faire sont indiqués progressivement.}

\section*{Objectifs de fin de séquence}

L'élève doit pouvoir :

\begin{itemize}
  \item identifier les deux grandeurs étudiées et préciser leurs unités ;
  \item reconnaître une situation de proportionnalité et expliquer son choix ;
  \item compléter un tableau de proportionnalité par une méthode adaptée ;
  \item utiliser un coefficient, un retour à l'unité, la multiplication ou division
        d'une colonne, l'additivité ou la soustraction ;
  \item résoudre des problèmes de prix, de recette, d'échelle ou d'agrandissement ;
  \item rappeler population, caractère, modalité, effectif, fréquence et pourcentage ;
  \item lire, contrôler et construire un diagramme circulaire ;
  \item passer d'un effectif ou d'un pourcentage à un angle ;
  \item exercer son esprit critique face à une représentation médiatique.
\end{itemize}

\section*{Déroulé réaliste}

\renewcommand{\arraystretch}{1.38}
\begin{tabularx}{\textwidth}{
|>{\centering\arraybackslash}p{1.05cm}
|>{\raggedright\arraybackslash}p{4cm}
|>{\raggedright\arraybackslash}X
|>{\raggedright\arraybackslash}p{3.1cm}|}
\hline
\textbf{Séance} & \textbf{Support principal} & \textbf{Enjeu didactique} & \textbf{Production}\\
\hline
1 & Activité 1 : classer des situations & Identifier les grandeurs et leurs unités ; ne pas supposer que toute situation à deux grandeurs est proportionnelle. & Tableau argumenté.\\
\hline
2 & Activité 2 : puzzle agrandi & Faire vivre le coefficient d'agrandissement comme opérateur commun à toutes les longueurs. & Puzzle et calculs.\\
\hline
3 & Activité 3 : compléter des tableaux & Comparer additivité, linéarité multiplicative, coefficient et retour à l'unité. & Tableaux complétés.\\
\hline
4 & Activité 4 : fruits, boucher et maquette & Choisir les grandeurs pertinentes ; résoudre dans les deux sens ; travailler les unités. & Interrogation 1 au début.\\
\hline
5 & Activité 5 : enquête statistique & Réactiver effectifs, fréquences et pourcentages ; vérifier les totaux. & Tableau statistique.\\
\hline
6 & Activité 6 : diagrammes et médias & Contrôler la somme des pourcentages et la cohérence entre nombres et secteurs. & Critique argumentée.\\
\hline
7 & Activité 7 : construire un diagramme & Relier effectif, pourcentage et angle par proportionnalité ; construire au rapporteur. & Interrogation 2 puis diagramme.\\
\hline
8 & Consolidation différenciée & Choix de méthode, analyse d'erreurs, tableaux et diagrammes mixtes. & Synthèse.\\
\hline
9 & Évaluation finale & Mobiliser reconnaissance, calcul, représentation et esprit critique. & Évaluation sur 20.\\
\hline
\end{tabularx}

\section*{Automatismes}

Rituel de cinq minutes, quatre à six questions :

\begin{itemize}
  \item séances 1 à 3 : reconnaître une situation et nommer les grandeurs ;
  \item séances 2 à 5 : compléter une colonne ou une ligne de tableau ;
  \item séances 4 à 7 : calculer une quatrième proportionnelle ou revenir à l'unité ;
  \item séances 5 à 8 : calculer fréquence, pourcentage ou angle ;
  \item après la séquence : une question de proportionnalité et une lecture de données par semaine.
\end{itemize}

\section*{Principes didactiques}

\begin{itemize}
  \item La reconnaissance de la proportionnalité précède toute technique.
  \item Plusieurs procédures correctes sont acceptées et comparées.
  \item Le produit en croix n'est pas présenté comme une recette automatique.
  \item Des situations non proportionnelles sont régulièrement proposées.
  \item Le tableau effectif--angle donne du sens à la construction du diagramme circulaire.
  \item Les diagrammes médiatiques sont contrôlés numériquement et visuellement.
\end{itemize}

\end{document}
